\section*{第 \ref{ch:b3:continuum-elasticity} 章——连续介质力学与弹性}

\begin{solution}{exo:b3:continuum-elasticity:1}
(a) $\unit{Pa} = \unit{N/m^2} = \unit{kg.m^{-1}.s^{-2}}$；应变是
无量纲数。(b) $\varepsilon = \sigma/E = \num{e-3}$。(c) 长度加倍
即 $\varepsilon = 1$，超出“小”的范围一千倍，
橡胶在此处的响应强烈非线性（起源于熵，性质不同），
胡克定律只是曲线的开端。(d) $w =
\sigma^2/2E$：钢 $(\num{2e8})^2/2(\num{2e11}) =
\qty{1e5}{J/m^3}$；橡胶 $(\num{e6})^2/2(\num{2e6}) =
\qty{2.5e5}{J/m^3}$——柔软材料在工作应力下单位体积
储存的能量\emph{更多}，这正是弹弓用橡胶而非
钢制的原因。
\end{solution}

\begin{solution}{exo:b3:continuum-elasticity:2}
(a) $\varepsilon_{ij} = \alpha\delta_{ij}$：各向同性膨胀，
$\delta V/V = 3\alpha$。(b) $\varepsilon_{xy} = \varepsilon_{yx} =
\gamma/2$，其余为零：纯剪切，体积不变。(c)
$\varepsilon_{ij} = 0$：刚性转动，无形变。(d)
$\operatorname{diag}(\varepsilon, -\nu\varepsilon, -\nu\varepsilon)$：
拉伸试验——沿 $x$ 方向拉伸，横向为泊松收缩。
\end{solution}

\begin{solution}{exo:b3:continuum-elasticity:3}
(a) 仅有 $\sigma_{zz}(z)$ 与重力作用时：$\partial_z\sigma_{zz} =
\rho g$（压缩取向下为正），$\sigma_{zz} = \rho
gh$，即底部的值。(b) $h_{\max} = \sigma_{\text{c}}/\rho g =
\num{2e8}/(2700 \times 9.81) \approx \qty{7.6}{km}$。(c) 珠穆朗玛峰正处于
其自身底部的压溃极限上——地球上的山脉
高度恰受岩石强度限制，不会更高。(d) 高为 $h$ 的锥体对其底部
仅施加 $\rho gh/3$（为柱体的三分之一：质量随
截面增大而增加），因此山形堆积体可以
站立到约为柱体三倍的高度。
\end{solution}

\begin{solution}{exo:b3:continuum-elasticity:4}
(a) $-\nu\varepsilon = \num{-3e-4}$。(b) $\delta V/V = (1 -
2\nu)\varepsilon = \num{4e-4}$。(c) $\nu = 1/2$：不可压缩
形变，即橡胶。(d) 当 $\nu > 1/2$ 时，体积模量 $K =
E/3(1 - 2\nu)$ 将为负：四面挤压会
\emph{增大}体积，材料会通过坍缩来
释放能量——任何稳定的固体都做不到。
\end{solution}

\begin{solution}{exo:b3:continuum-elasticity:5}
(a) 对单位长度的半圆柱，压力以 $p
\times 2R$（投影面积）推动，两个切开的壁面则以
$2\sigma_\theta t$ 的力反向拉回：$\sigma_\theta = pR/t$。(b) 在横截面上，
$p\pi R^2 = \sigma_z\,2\pi Rt$：$\sigma_z = pR/2t$，即一半。
香肠沿\emph{纵向}裂开：较大的环向应力会
沿轴向撕裂外皮。(c) $\sigma_\theta = \num{2e7} \times
0.09/0.005 = \qty{360}{MPa}$：较屈服强度低 $2$ 倍——这正是
圆柱体要经过强度检验并接受检查的原因。(d) 球体在各个方向承受
$pR/2t$：半球形端部使应力减半，
并避免平端在边角处集中应力。
\end{solution}

\begin{solution}{exo:b3:continuum-elasticity:6}
(a) $G = E/2(1 + \nu) = \qty{1.0}{MPa}$。(b) $\sigma_{xy} = F/S =
150/\num{e-2} = \qty{15}{kPa}$；$\gamma = \sigma_{xy}/G = 0.015$；
位移 $\gamma h = \qty{0.3}{mm}$。(c) 橡胶网络改变
\emph{形状}只需使其分子链重新取向（较容易），而改变
\emph{体积}则意味着使分子链更紧密地堆积，难度堪比压缩
液体：$G \ll K$，因此 $\nu \to 1/2$。(d) 剪切方向柔软：
隔震支座允许地面在其下方水平晃动，同时传递的
力很小，但在压缩方向上仍足够刚硬，以承托
建筑物的重量。
\end{solution}

\begin{solution}{exo:b3:continuum-elasticity:7}
(a) $k = ES/L = \num{1.2e9} \times \num{8e-5}/20 = \qty{4.8}{kN/m}$。
(b) $mg(h + x) = \tfrac12 kx^2$：$2400x^2 - 785x - 3140 = 0$，$x =
\qty{1.3}{m}$。(c) $F = kx \approx \qty{6.3}{kN}$；减速度 $F/m
\approx \qty{79}{m/s^2} \approx 8g$——绳索之所以
要刻意做得有弹性。(d) 被吸收的能量中有一部分
用于使纤维断裂和塑性重排：绳索现处于一条
退化的应力-应变曲线上，更硬而更弱，下一次坠落
在两方面都会更加严峻。
\end{solution}

\begin{solution}{exo:b3:continuum-elasticity:8}
(a) 半径为 $r$ 的管顶部转过的弧长为 $r\theta$，
在长度 $L$ 上：$\gamma = r\theta/L$。(b) $\Gamma = \int_0^a r\,(G r\theta/L)\,
2\pi r\,\dd r = (\pi Ga^4/2L)\,\theta$。(c) 增大为 $2^4 = 16$ 倍。(d) $C =
\pi \times \num{4e10} \times (\num{2.5e-5})^4/2 =
\qty{2.5e-8}{N.m/rad}$；$T = 2\pi\sqrt{I/C} = 2\pi\sqrt{\num{5e-4}/
\num{2.5e-8}} \approx \qty{900}{s}$——每次摆动约需四分之一小时，
灵敏度足以感受铅球的引力。
\end{solution}

\begin{solution}{exo:b3:continuum-elasticity:9}
(a) $\mu = E/2(1+\nu) = \qty{78}{GPa}$，$\lambda = E\nu/(1+\nu)(1-2\nu)
= \qty{107}{GPa}$：$c_{\text{P}} = \sqrt{262\times10^9/7850} =
\qty{5.8}{km/s}$，$c_{\text{S}} = \qty{3.1}{km/s}$。(b) 细棒
给出 $\sqrt{E/\rho} = \qty{5.0}{km/s}$：细棒中侧面可以自由鼓出
（泊松释放），使响应变软；大块介质中
周围物质不允许这种鼓出。(c) $c_{\text{S}} = 0$；
$c_{\text{P}} = \sqrt{\num{2.2e9}/1000} = \qty{1.5}{km/s}$——即
水中的声速。(d) P 波沿射线推动地面
（深源时近乎竖直），S 波则垂直于它：
地震仪的三个分量可分辨两种偏振，进而由 P/S 延迟
推定震源距离。
\end{solution}

\begin{solution}{exo:b3:continuum-elasticity:10}
(a) 半径为 $R + y$ 的弧长为 $(R + y)\alpha$，而
中性面处为 $R\alpha$：$\varepsilon = y/R$，$\sigma =
Ey/R$。(b) $M = \int y\sigma\,\dd S = (E/R)\int y^2\,\dd S = EI/R$；
对矩形截面 $I = bh^3/12$。(c) 平放：$I = bh^3/12$；立放：
$hb^3/12$——比值 $(b/h)^2 = 25$（$b/h = 5$）：同一木材，刚度
提高为二十五倍，纯粹由几何决定。(d) $I = 0.50 \times
0.035^3/12 = \qty{1.8e-6}{m^4}$，$EI = \qty{2.1e4}{N.m^2}$；$\delta =
FL^3/3EI = 736 \times 1.8^3/(3 \times \num{2.1e4}) \approx
\qty{7}{cm}$：恰好就是跳水板那种令人愉快的弹性。
\end{solution}

\begin{solution}{exo:b3:continuum-elasticity:11}
(a) 在弯曲构型中，轴向载荷 $P$ 的力臂为
$y(x)$，截面位于 $x$ 处：$M = -Py$，而 $M = EIy''$ 给出
$EIy'' = -Py$。(b) $y = A\sin(x\sqrt{P/EI})$，由 $y(L) = 0$ 知非零的
$A$ 首次出现在 $\sqrt{P/EI}\,L = \pi$ 处，即 $P_{\text{c}} =
\pi^2EI/L^2$；低于此值时直柱为唯一解，高于此值时
弯曲无需额外的力。(c) $I = 0.03 \times 0.0015^3/12 =
\qty{8.4e-12}{m^4}$：$P_{\text{c}} = \pi^2 \times 12\times10^9 \times
\num{8.4e-12}/1^2 \approx \qty{1}{N}$——相当于一只苹果的重量：一把
米尺在手指轻压之下即会屈曲。(d) $I$ 随材料远离轴线而增大：
空心管把全部截面置于较大的 $y$ 处，
在材料重量一定时使 $I$（进而 $P_{\text{c}}$）达到最大——
骨骼、竹子和自行车车架正是如此设计的。
\end{solution}

\begin{solution}{exo:b3:continuum-elasticity:12}
(a) 每个面积为 $a^2$ 的区域对应一条链；应变 $\varepsilon$ 拉伸时
每根弹簧伸长 $\delta = \varepsilon a$，每条链的力为 $\kappa
\varepsilon a$，应力为 $\kappa\varepsilon/a$：$E =
\kappa/a$。(b) $\rho = m/a^3$，于是 $E/\rho = \kappa a^2/m$，$c =
a\sqrt{\kappa/m}$。(c) $c = \num{2.5e-10}\sqrt{50/\num{e-25}} \approx
\qty{5.6}{km/s}$。(d) 恰与金属实测的几千米每秒相吻合：一切固体的
刚度——以及每一次地震波的速度——
都取决于化学键的刚度。
\end{solution}

\begin{solution}{pb:b3:continuum-elasticity:1}
\textbf{1.} $\mu = E/2(1+\nu) = \qty{30}{GPa}$；$\lambda =
E\nu/(1+\nu)(1-2\nu) = \qty{30}{GPa}$：当 $\nu = 1/4$ 时，$\lambda =
\mu$。
\textbf{2.} $c_{\text{P}} = \sqrt{3\mu/\rho} = \sqrt{\num{9e10}/2700}
= \qty{5.8}{km/s}$；$c_{\text{S}} = \sqrt{\mu/\rho} =
\qty{3.3}{km/s}$；比值为 $\sqrt3$。
\textbf{3.} $\varepsilon \sim 2/\num{5e4} = \num{4e-5}$；$\sigma =
E\varepsilon \approx \qty{3}{MPa}$——真实地震中测得的“应力降”
确实只有几兆帕。
\textbf{4.} $w = \tfrac12 E\varepsilon^2 \approx \qty{60}{J/m^3}$；
乘以 $\num{5e13}{}\,\unit{m^3}$：$\sim\qty{3e15}{J}$，约相当于
四分之三兆吨 TNT——一次大地震。
\textbf{5.} 从下方来的 P 波使地面竖直运动——建筑物本就按承受竖直荷载设计；
S 波则使地面\emph{水平}运动，
而建筑物在这一方向上最为薄弱。
\textbf{6.} 只有 P 波：它是压缩波；S 波是纯剪切，
体积保持不变。
\textbf{7.} $t_{\text{P}} = d/c_{\text{P}}$，$t_{\text{S}} =
d/c_{\text{S}}$：$\Delta t = d(1/c_{\text{S}} - 1/c_{\text{P}})$，
由此求逆即得题中所述。
\textbf{8.} $1/(1/3.33 - 1/5.77)\,\unit{km/s} \approx
\qty{7.9}{km}$ 每秒延迟——即地震学家所说的
“每秒八公里”。
\textbf{9.} $d \approx 7.9 \times 28 \approx \qty{220}{km}$。
\textbf{10.} 一般需要三个台站：每个台站确定一个可能震中的圆，
两圆相交于两点，第三个圆
决定其中的一点。
\textbf{11.} $t_{\text{P}} = 80/5.8 = \qty{14}{s}$，$t_{\text{S}} =
80/3.3 = \qty{24}{s}$：自震动开始到破坏性摇晃之间，
约有 \qty{10}{s} 的预警时间。
\textbf{12.} S 波在 $200/3.3 = \qty{60}{s}$ 时到达。P 波
于 \qty{3.5}{s} 时到达 \qty{20}{km} 处；无线电信号几乎瞬时到达，
因此该城市可获得近一分钟的预警——
现有系统可达数十秒。
\textbf{13.} $2R_{\text{E}}/c \approx 12742/10 \approx \qty{1270}{s}$，
约二十分钟。
\textbf{14.} 弹性波竟能穿过它，说明深部地球并非通体熔融——
它能传递波动，并且（S 波穿过地幔时）
甚至像固体一样发生剪切。
\textbf{15.} 两个半径以及它们之间的夹角 $\Delta$：弦长为
$2R_{\text{E}}\sin(\Delta/2)$。
\textbf{16.} $2 \times 6371 \times \sin30^\circ = \qty{6371}{km}$；
$t \approx \qty{640}{s} \approx \qty{11}{min}$。
\textbf{17.} 超过 $103^\circ$ 后，每条射线都必须穿过地球深处的
中心区域；若该区域为\emph{流体}（$\mu = 0$），则它不传播
剪切波：S 波在那里消失。流体正是
不能抵抗剪切的物态。
\textbf{18.} 流体仍能传播压缩：P 波穿过地核，
但速度较慢——因此它们在其边界处发生明显折射，
弯曲的射线留下一个环形阴影，而非清晰的截断。
\textbf{19.} 液态地核位于地球中心。
\textbf{20.} 弦距地心最近的距离为
$R_{\text{E}}\cos(\Delta/2)$；掠过地核的条件为
$R_{\text{E}}\cos(\Delta^*/2) = R_{\text{c}}$，即 $\Delta^* =
2\arccos(R_{\text{c}}/R_{\text{E}})$。
\textbf{21.} $R_{\text{c}} = 6371\cos(51.5^\circ) \approx
\qty{3970}{km}$。
\textbf{22.} 偏大 $+14\%$。速度随深度增加，真实射线因而
向地表弯回：最深点恰好触及地核的射线，其出射角
\emph{小于}穿过同一深度的直线弦的角度——
因此由观测到的 $103^\circ$ 出发，
用直线射线反演把切点定得过浅，即
高估了地核。
\textbf{23.} 阴影区内存在微弱的 P 波能量，说明液态地核内部有
某种结构使射线折回：存在一个独立的\emph{内核}
（莱曼，1936）。
\textbf{24.} 剪切波只能在固体中传播：一个 P 波在内核内
转换为剪切波、再转换回来（即 PKJKP 震相），
便可证明其为固体——这一
长期追寻的探测，已被公认为证据。
\textbf{25.} 岩石中的两个速度（$5.8$ 与 \qty{3.3}{km/s}）、
$103^\circ$ 以外缺失的一种波，以及弦的几何关系，共同给出
液态地核约 $\approx\qty{4000}{km}$——与现代值
\qty{3480}{km} 相差在 $15\%$ 以内：胡克定律在行星尺度上的应用。
\end{solution}
